% =============================================================================
%  Chapter 3 -- Methodology
% =============================================================================
\chapter{Methodology}
\label{sec:methodology}

\section{Mathematical notation}
\label{sec:notation}

% NOTATION. Fix a convention in this section and never deviate from it:
%   scalars       x        lowercase italic
%   vectors       \bm{x}   lowercase bold        (use \bm, not \mathbf)
%   matrices      \bm{X}   uppercase bold
%   sets/spaces   \mathcal{H}, \mathbb{R}
%
% \bm{} is loaded by the class and bolds maths correctly, including Greek
% letters and symbols; \mathbf{} silently fails on \alpha and friends.
%
% Define every symbol at first use, and keep text/notations.tex in step: that
% front-matter list is written by hand and nothing checks it against the text.

Throughout, scalars are written $x$, vectors $\bm{x}$, matrices $\bm{X}$, and
the hypothesis space $\mathcal{H} \subseteq \mathbb{R}^{d}$.

\section{Equations}
\label{sec:equations}

% EQUATIONS. Three rules that cover almost every case:
%
%  1. Punctuate them. A displayed equation is part of a sentence and takes the
%     comma or period the sentence requires, as below.
%  2. Number only what you reference. Use equation* (or \[ ... \]) otherwise --
%     a column of unreferenced numbers is noise.
%  3. Never use $$ ... $$. It is plain TeX, and it breaks spacing above and
%     below the display.
%
% Reference with \eqref{eq:...}, which supplies the parentheses: "(3.1)".

The cross-entropy loss over $N$ observations and $K$ classes is defined as
\parencite{goodfellow2016deep}
\begin{equation}
  \mathcal{L}(\bm{\theta}) = -\frac{1}{N} \sum_{i=1}^{N} \sum_{k=1}^{K}
    y_{ik} \log \hat{y}_{ik}(\bm{\theta}),
  \label{eq:cross_entropy}
\end{equation}
where $\hat{y}_{ik}$ denotes the predicted probability that observation $i$
belongs to class $k$. Minimising \eqref{eq:cross_entropy} with respect to
$\bm{\theta}$ yields the fitted model.

% An unnumbered display, for an aside you will never point back to:
\begin{equation*}
  z_t = \frac{y_t - \mu}{\sigma}.
\end{equation*}

\section{Definitions and theorems}
\label{sec:theorems}

% Four environments are predefined by the class: definition, theorem,
% proposition, lemma. They are numbered per chapter (Definition 3.1, 3.2, ...).
% Add more with \newtheorem in preamble.tex if you need them.

\begin{definition}[Named definition]
\label{def:example}
State the object being defined precisely, and give it a name in the brackets so
it can be referred to by name rather than only by number.
\end{definition}

\begin{theorem}
\label{thm:example}
State the claim. If you prove it, the proof follows in a \texttt{proof}
environment, which ends with an automatic QED square.
\end{theorem}

\begin{proof}
The actual proof.
\end{proof}

\section{Model architecture}
\label{sec:architecture}

% STYLE: describe the model in enough detail that someone could rebuild it:
% the layers and their sizes, the activation, what is trainable and what is
% frozen, and where the numbers came from. A hyperparameter that was tuned
% should say so; one inherited from a paper should cite it.
%
% A diagram earns its place here more than almost anywhere else in the thesis,
% because an architecture is far quicker to see than to read. See
% \Cref{sec:figures} for the conventions.

Placeholder description of the model, also shown in
\Cref{fig:multihead_attention}.

\begin{figure}
\centering
% =============================================================================
%  Multi-head attention.
%
%  Adapted from NNTikZ, Copyright (c) 2024 Fraser Love, MIT Licence:
%  https://github.com/fraserlove/nntikz -- re-coloured with this template's
%  palette and reduced to a bare fragment.
%
%  Deliberately restrained: a single hue (tolblue) marks the blocks, everything
%  else is black. The network figures in Chapter 2 use several colours to show
%  the palette off; a figure whose parts carry no categorical distinction does
%  not need more than one.
%
%  Needs the TikZ libraries chains, quotes, decorations.pathreplacing and
%  ext.paths.ortho, all loaded by aml-thesis.cls.
% =============================================================================

% Machinery for the stacked ("shadowed") blocks that represent the H heads.
\makeatletter
\tikzset{phantom node/.code=\tikz@addoption{\expandafter\let\csname pgf@sh@boxes@\tikz@shape\endcsname\pgfutil@empty}}
\makeatother
\tikzset{
  shadowed node xshift/.initial=1.5ex,
  shadowed node yshift/.initial=1ex,
  shadowed node list/.initial={2, 1},
  pics/shadowed node/.default=\pgfkeysvalueof{/tikz/shadowed node list},
  shadowed node/.pic={
    \foreach[expand list] \elem in {#1} {
      \scoped[transparency group, shadowed node calculation={\elem}]
        \node[style/.expand once=\tikzpictextoptions, phantom node,
              xshift={\elem*\pgfkeysvalueof{/tikz/shadowed node xshift}},
              yshift={\elem*\pgfkeysvalueof{/tikz/shadowed node yshift}}]
              (-\elem) {\tikzpictext};
    }
    \node[alias=-0, style/.expand once=\tikzpictextoptions] () {\tikzpictext};
  },
  set shadowed node calculation parameter/.style={
    shadowed node calculation/.style={opacity={(#1-##1+1)/(#1+1)}}
  },
  set shadowed node calculation parameter=2,
  overshoot line to/.style={
    to path={($(\tikztostart)!-(#1)!(\tikztotarget)$)--($(\tikztotarget)!-(#1)!(\tikztostart)$)\tikztonodes}
  },
  edges have transparency group/.style={
    execute at begin to={\scope[transparency group,#1]},
    execute at end to=\endscope
  }
}

\begin{tikzpicture}[
    >=LaTeX,
    thick,
    x=2cm,
    node distance=1.75em,
    cell/.style={rectangle, rounded corners=2mm, minimum height=2em,
                 minimum width=2.5em, draw=tolblue, fill=tolblue!10, thick},
    arrow/.style={-latex, very thick, draw=black}
]

    \node foreach[count=\i] \t in {$\bm{Q}$, $\bm{K}$, $\bm{V}$} (VKQ-\i) at (\i, 0) {\t}
    pic foreach \i in {1, 2, 3}
        ["Linear" cell, above=of VKQ-\i] (Linear-\i) {shadowed node}
    pic["Scaled Dot-Product Attention" cell, above = of Linear-2]
        (Attention) {shadowed node}
    node[above = of Attention, cell] (Concat) {Concat}
    node[above = of Concat, cell] (Linear) {Linear}
    node[above = of Linear] (MHAtt) {Multi-Head Attention}
    [arrow, ortho/install shortcuts]
    foreach \i in {1, 2, 3}{
        (VKQ-\i) edge coordinate [pos=.2] (@) (Linear-\i)
        (Linear-\i) edge[|*] (Attention)
        foreach \j in {2, 1}{
        [
            edges have transparency group={shadowed node calculation=\j},
            behind path
        ]
        (@) edge[out=90, in=-90] (Linear-\i-\j)
        [in front of path]
        (Linear-\i-\j) edge[path only, |*] coordinate (@@) (Attention-\j)
                        edge[-] (@@)
        [behind path] (@@) edge[|*] (Attention-\j)
        }
    }
    foreach \i in {0, 1, 2}{
        [edges have transparency group={shadowed node calculation=\i}]
        (Attention-\i) edge[|*] (Concat)
    }
    (Concat) edge (Linear)
    (Linear) edge (MHAtt);

    % Brace marking the H parallel heads.
    \draw[
        arrows={[arc=135]}, arrows={Hooks[left]-Hooks[right]},
        s/.style={shift={(3mm,-3.5mm)}}
    ]
    ([s]Attention-2.east) to[overshoot line to=1mm] coordinate(@) ([s]Attention.east);
    \path node[right=3mm] at (@) {$H$};
\end{tikzpicture}

\caption[Multi-head attention]{Multi-head attention
\parencite{vaswani2017attention}: $H$ heads each project $\bm{Q}$, $\bm{K}$ and
$\bm{V}$ through their own linear map and attend in parallel, and their outputs
are concatenated and projected once more. Adapted from
\textcite{love2024nntikz}.}
\label{fig:multihead_attention}
\end{figure}

\section{Algorithms and code}
\label{sec:algorithms}

% CODE. Use listings for real source code you want reproduced verbatim, and a
% dedicated algorithm package (algorithm2e, algpseudocode) for pseudo-code.
% Keep listings short -- more than ~20 lines belongs in the appendix or a
% repository, with a URL in a footnote.

\begin{lstlisting}[language=Python,caption={Keep listings short and captioned.},label={alg:example}]
def train(model, loader, optimiser):
    for batch in loader:
        loss = model.loss(batch)
        loss.backward()
        optimiser.step()
    return model
\end{lstlisting}
